\documentclass[10pt, a4paper]{article}
% ===== Essential Packages =====
\usepackage[utf8]{inputenc}       % UTF-8 encoding
\usepackage[T1]{fontenc}          % Better font encoding
\usepackage{lmodern}              % Modern Latin Modern font
\usepackage{ctex}                 % Chinese support (must be loaded after fontenc)

% ===== Math & Symbols =====
\usepackage{amsmath}              % Advanced math
\usepackage{amssymb}              % Additional math symbols
\usepackage{mathrsfs}             % Ralph Smith's Formal Script
\usepackage{braket}               % Dirac bra-ket notation

% ===== Graphics & Tables =====
\usepackage{graphicx}             % Image support
\usepackage{float}                % Better float control
\usepackage{wrapfig}              % Wrapped figures
\usepackage{tikz}                 % Vector graphics
\usepackage[table]{xcolor}        % Colored tables
\usepackage{array}                % Enhanced tables
\usepackage{multirow}             % Multirow tables
\usepackage{tabularx}             % Adjustable-width tables
\usepackage{booktabs}             % Professional table lines
\usepackage{siunitx}              % SI units in tables
\usepackage{caption}              % Better captions
\usepackage{adjustbox}            % Box resizing

% ===== Code Listings =====
\usepackage{listings}             % Code blocks
\usepackage{tcolorbox}            % Colored boxes
\tcbuselibrary{listings, breakable} % Listings in tcolorbox

% ===== Misc =====
\usepackage{lipsum}               % Dummy text
\usepackage{footmisc}             % Better footnotes
\usepackage[margin=1in]{geometry} % Page margins
\title{焦距的测定}
\author{张世航}
\begin{document}
	\maketitle
	\section{自准法测凸透镜焦距}
	\begin{tabular}{|c|c|c|c|c|c|c|}
		\hline
		$x_1(cm)$& \multicolumn{6}{|c|}{$x_2(cm)$}\\
		\hline
		$10.00$& 24.70 & 24.80 & 24.90 & 25.00 & 25.10 & 25.20   \\
		\hline
		$10.00$& 25.20 & 25.13 & 25.03 & 24.95 & 24.85 & 24.83   \\
		\hline
	\end{tabular}
	
	\section{位移法测凸透镜焦距}
	\begin{tabular}{|c|c|c|c|c|c|c|c|}
		\hline
		$N$& $A_1(cm)$ & $A_2(cm)$ & $A=A_1-A_2(cm)$ & $L_1(cm)$ & $L_2(cm)$ & $L=L_1-L_2(cm)$ &$f=\frac{A^2-L^2}{4A}(cm)$  \\
		\hline
		1& 10.00 & 100.00 & 90.00 & 28.62 & 81.62 & 53.00 & 14.67 \\
		\hline
		2& 10.00 & 100.00 & 90.00 & 28.60 & 81.64 & 52.96 & 14.70 \\
		\hline
		3& 10.00 & 100.00 & 90.00 & 28.62 & 81.58 & 52.96 & 14.70 \\
		\hline
		4& 10.00 & 100.00 & 90.00 & 28.65 & 81.61 & 52.96 & 14.70 \\
		\hline
		5& 10.00 & 100.00 & 90.00 & 28.61 & 81.61 & 53.00 & 14.67 \\
		\hline
		6& 10.00 & 100.00 & 90.00 & 28.58 & 81.65 & 53.01 & 14.69 \\
		\hline
		平均值&10.00&100.00&90.00&28.613&81.618&52.982&14.688\\
		\hline
	\end{tabular}
	
	\section{误差分析}
	\subsection{自准法}
	利用A类误差公式\[u_A = \sqrt{\frac{\sum_{i=1}^n(x_i-\bar{x})^2}{n(n-1)}}t_p\]
	我们取\[t_p = 0.683\]
	和B类误差公式
	\[u_B = \frac{\Delta_{\text{仪器}}}{C}\]
	我们取\[C = \sqrt{3}\]
	我们可以得到
	\[u_{x_1} = 0.06612cm\]
	\[u_{x_2} = 0.05774cm\]
		根据
	\[U_r = \sqrt{\sum_{i=1}^n(\frac{\partial lnf}{\partial x_i})^2u_c^2(x_i)}\]
	\[U = \bar{f}U_r = 0.08778cm\]
	则\[f = (14.974\pm0.08778)cm\]
	\subsection{位移法}
	利用A类误差公式\[u_A = \sqrt{\frac{\sum_{i=1}^n(x_i-\bar{x})^2}{n(n-1)}}t_p\]
	我们取\[t_p = 0.683\]
	和B类误差公式
	\[u_B = \frac{\Delta_{\text{仪器}}}{C}\]
	我们取\[C = \sqrt{3}\]
	我们可以得到
	\[u_{A_1} = 0.05810194200235001cm\]
	\[u_{A_2} = 0.058148757820878cm\]
	\[u_{L_1} = 0.05773502691896258cm\]
	\[u_{L_2} = 0.05773502691896258cm\]
	根据
	\[U_r = \sqrt{\sum_{i=1}^n(\frac{\partial lnf}{\partial x_i})^2u_c^2(x_i)}\]
	\[U = \bar{f}U_r\]
	\[\frac{\partial lnf}{\partial A_1}= \frac{1}{f}(\frac{1}{4}+\frac{L^2}{4A^2})=0.022919cm^{-1}\]
	\[\frac{\partial lnf}{\partial A_2}= -\frac{1}{f}(\frac{1}{4}+\frac{L^2}{4A^2})=0.022919cm^{-1}\]
	\[\frac{\partial lnf}{\partial L_1} = -\frac{1}{f}\frac{L}{2A}=0.02004cm^{-1}\]
	\[\frac{\partial lnf}{\partial L_2} = \frac{1}{f}\frac{L}{2A}=0.02004cm^{-1}\]
	\[U = fU_r = 0.03651157204495233cm\]
	则\[f = (14.688\pm0.0365)cm\]
	
	
    $\int_0^\infty(-2r)e^{-4\alpha r}dr$
\end{document}